\chapter*{Abstract}
\addcontentsline{toc}{chapter}{Abstract}

This study presents a determinant-based workflow for predicting antimicrobial resistance (AMR) phenotypes from curated genomic evidence. The intended target is not one bacterial species, but the broader genotype-to-phenotype prediction problem for organisms with sufficiently rich paired genotype and phenotype data. \textit{Salmonella enterica} is used as the case study because it is well represented in surveillance resources and provides enough data for a credible antibiotic-specific evaluation. Rather than using genome-wide features such as k-mers or pan-genome encodings, the study focuses on biologically interpretable AMR determinants derived from AMRFinderPlus and the NCBI MicroBIGG-E resource.

The workflow aligns phenotype and genotype records through shared \texttt{BioSample} identifiers, cleans AST labels into a stable binary resistant/susceptible target, deduplicates and normalizes determinant annotations, and then constructs antibiotic-specific model-input tables under three genotype-to-phenotype connection scopes: \textit{strict}, \textit{broad}, and \textit{all}. Antibiotic-specific XGBoost models are trained under a consistent cross-validation protocol, and direct rule-based baselines are included under constrained scopes to provide a deterministic biological reference.

The results show a clear pattern. The \textit{all} scope provides the strongest overall predictive performance and removes the zero-feature problem completely. The \textit{broad} scope is the strongest biologically constrained alternative and performs especially well when broader within-class resistance mechanisms must be captured. The \textit{strict} scope remains useful as a narrow biological reference, but it is too sparse to serve as a universal default. Overall, the study shows that determinant-based AMR prediction is both practical and interpretable, and that the genotype-to-phenotype connection rule is itself a major scientific variable rather than a minor implementation detail.
